% ══════════════════════════════════════════════════════════════════════════════
%  LESSON 1 — WHAT IS THE WATER CYCLE?                                  (Week 1 · P2)
% ══════════════════════════════════════════════════════════════════════════════
\setcounter{section}{0}
\sectiondiv{What is the Water Cycle?}{Evaporation, condensation and precipitation \textbullet{} Collection and transpiration \textbullet{} Why fresh water is limited}\label{sec:1}
\lessoncodes{VC2S8U09 \textbullet{} VC2S8U05}

\vocabhead
\begin{ctab}
\begin{tabular}{|SW{3.4cm}|W{7.0cm}|}
  \hline\hd & \thead{Word} & \thead{Definition (copy from the slide)} \\ \hline
  & \scaf{\textbf{Water cycle}} & \scaf{The continuous movement of water between the land, the oceans and the air.} \\ \hline
  & \textbf{Evaporation} & \\ \hline
  & \textbf{Condensation} & \\ \hline
  & \textbf{Precipitation} & \\ \hline
  & \textbf{Transpiration} & \\ \hline
  & \textbf{Renewable resource} & \\ \hline
\end{tabular}
\end{ctab}

% ─── PART A ───────────────────────────────────────────────────────────────────
\lessonpart{A}{Reading and comprehension}{Teal}\label{read:1}

\begin{readingbox}{Reading --- The Same Water, Again and Again}
The water you drink today may once have fallen as rain on a dinosaur. Earth does not make new water; it recycles the same water over and over in the \textbf{water cycle}.

The Sun drives the cycle. Its energy warms the oceans, lakes and rivers, and water at the surface turns into water vapour, a gas. This is \textbf{evaporation}. Plants add more water vapour through tiny holes in their leaves: \textbf{transpiration}.

The warm, moist air rises. High in the sky it cools, and the water vapour turns back into tiny liquid droplets that gather into clouds. This is \textbf{condensation}. When the droplets join and grow too heavy, they fall as rain, hail or snow: \textbf{precipitation}.

Rain that lands on the ground runs into creeks and rivers, soaks into the soil to become groundwater, or collects in lakes, dams and the sea. Then the cycle starts again.

Because the cycle never stops, water is a \textbf{renewable resource}. Yet fresh water is still scarce. About 97\% of Earth's water is salty sea water, and most of the fresh water is frozen in ice sheets. Less than 1\% of all the water on Earth is fresh water we can easily reach.
\end{readingbox}

\rcreset
\rcq{What drives the water cycle?}{2.5cm}
\rcq{What is the difference between evaporation and transpiration?}{2.5cm}
\rcq{Explain how clouds form.}{2.5cm}
\rcq{Name three forms of precipitation.}{2.5cm}
\rcq{Give three places rain goes after it reaches the ground.}{2.5cm}
\rcq{If water is a renewable resource, why is fresh water still limited?}{3cm}

% ─── PART B ───────────────────────────────────────────────────────────────────
\lessonpart{B}{Theory and examples}{Blue}\label{theory:1}

\subsection{The Processes of the Water Cycle}\label{sub:1.1}

\begin{center}
\begin{tikzpicture}[scale=0.92, font=\small]
  \fill[black!12] (6.2,0) rectangle (11,0.9);
  \draw[Ink, line width=0.8pt] (6.2,0.9) -- (11,0.9);
  \draw[Ink, line width=1pt, fill=black!5] (0,0) -- (2.4,3.0) -- (4.6,0.9) -- (6.2,0.9) -- (6.2,0) -- cycle;
  \draw[Ink, line width=1pt] (4.2,0.9) -- (4.2,1.9); \draw[Ink, line width=1pt, fill=white] (4.2,2.3) circle (0.45);
  \draw[Ink, line width=0.8pt, densely dashed] (1.4,0.6) .. controls (3,0.3) and (5,0.5) .. (6.4,0.5);
  \fill[Ink] (2.2,5.0) circle (0.5); \fill[Ink] (2.9,5.3) circle (0.6); \fill[Ink] (3.6,5.0) circle (0.45); \fill[Ink] (2.2,4.6) rectangle (3.6,5.0);
  \fill[Ink] (8.2,5.0) circle (0.45); \fill[Ink] (8.8,5.25) circle (0.55); \fill[Ink] (9.4,5.0) circle (0.4); \fill[Ink] (8.2,4.6) rectangle (9.4,5.0);
  \foreach \x in {2.3,2.8,3.3}{\draw[Ink, line width=1pt] (\x,4.3) -- (\x-0.2,3.6);}
  \draw[Ink, line width=1pt] (10.6,6.0) circle (0.45);
  \foreach \x in {7.4,8.6,9.8}{\draw[Ink, line width=1.2pt, -{Stealth[length=5pt]}, densely dotted] (\x,1.0) -- (\x,3.8);}
  \draw[Ink, line width=1.2pt, -{Stealth[length=5pt]}, densely dotted] (4.2,2.8) -- (4.6,4.3);
  \draw[Ink, line width=1.2pt, -{Stealth[length=5pt]}] (7.9,5.0) -- (4.0,5.1);
  \node[draw=Ink, circle, inner sep=1.5pt, fill=white, font=\footnotesize\bfseries] at (8.9,2.4) {1};
  \node[draw=Ink, circle, inner sep=1.5pt, fill=white, font=\footnotesize\bfseries] at (4.95,3.6) {2};
  \node[draw=Ink, circle, inner sep=1.5pt, fill=white, font=\footnotesize\bfseries] at (8.8,4.3) {3};
  \node[draw=Ink, circle, inner sep=1.5pt, fill=white, font=\footnotesize\bfseries] at (1.95,3.9) {4};
  \node[draw=Ink, circle, inner sep=1.5pt, fill=white, font=\footnotesize\bfseries] at (5.2,1.35) {5};
  \node[draw=Ink, circle, inner sep=1.5pt, fill=white, font=\footnotesize\bfseries] at (3.2,0.35) {6};
  \node[font=\scriptsize] at (8.6,0.45) {sea};
  \node[font=\scriptsize] at (10.6,6.65) {Sun};
\end{tikzpicture}
\figcap{The water cycle. Each number marks one process.}
\end{center}

\begin{skillbox}{The processes}
\textbf{Evaporation:} liquid water turns into water vapour, using energy from the Sun. \quad \textbf{Transpiration:} plants release water vapour from their leaves. \quad \textbf{Condensation:} water vapour cools and turns into liquid droplets (clouds). \quad \textbf{Precipitation:} water falls from clouds. \quad \textbf{Collection:} water runs off into rivers and the sea. \quad \textbf{Infiltration:} water soaks into the ground (groundwater).
\end{skillbox}

\begin{workedex}{Processes as changes of state}
Which change of state happens in evaporation and in condensation?

\smallskip
\textbf{Evaporation:} liquid $\rightarrow$ gas (water $\rightarrow$ water vapour). It needs energy, from the Sun.\\
\textbf{Condensation:} gas $\rightarrow$ liquid (water vapour $\rightarrow$ droplets). It happens when the vapour cools.
\end{workedex}

\begin{wedo}{Processes as changes of state}
Snow falls on a mountain, and in spring it melts into a creek. Name the process in the water cycle and the change of state for each step.
\answerbox{4cm}
\end{wedo}

\subsection{Fresh Water is Limited}\label{sub:1.2}

\begin{workedex}{Earth's water as a bucket}
Imagine all of Earth's water is a 10-litre bucket. About 97\% is salty: $0.97 \times 10 = 9.7$ litres. The fresh water is only about 0.3 litres, and most of that is frozen. The fresh water we can easily reach is less than one cupful.
\end{workedex}

\begin{wedo}{Earth's water as a bucket}
Use the bucket to explain to a friend why we should not waste fresh water, even though the water cycle keeps renewing it.
\answerbox{4cm}
\end{wedo}

% ─── PART C ───────────────────────────────────────────────────────────────────
\lessonpart{C}{Practice}{Amber}\label{prac:1}

\task Name the process at each number on a water cycle diagram. The processes are: evaporation from the sea (1), transpiration from a tree (2), condensation into a cloud (3), precipitation from a cloud (4), collection as water runs off into the sea (5), and infiltration into the ground (6). Write the process for each number, and the change of state if there is one.
\begin{ctab}
\begin{tabular}{|SZ{1.2cm}|W{3.8cm}|W{4.6cm}|}
  \hline\hd & \thead{No.} & \thead{Process} & \thead{Change of state} \\ \hline
  & \scaf{1} & \scaf{Evaporation} & \scaf{Liquid to gas} \\ \hline
  & 2 & & \\ \hline
  & 3 & & \\ \hline
  & 4 & & \\ \hline
  & 5 & & \\ \hline
  & 6 & & \\ \hline
\end{tabular}
\end{ctab}

\needspace{10\baselineskip}
\task Wet washing dries faster on a hot, windy day. Explain why, using the word \emph{evaporation}.
\answerbox{5cm}

\needspace{10\baselineskip}
\task A cold can of drink gets wet on the outside on a warm day. Explain where the water comes from, using the word \emph{condensation}.
\answerbox{5cm}

\needspace{10\baselineskip}
\task Explain why water is called a renewable resource, but fresh water is still limited.
\answerbox{5cm}

% ─── PART D ───────────────────────────────────────────────────────────────────
\lessonpart{D}{Above level}{Violet}\label{ext:1}

\begin{yr8box}
These questions go beyond the Year 7 standard.

\medskip
\begin{keepext}
\exttask Clearing a large forest can reduce the rainfall in the area nearby. Use the water cycle to suggest why.
\answerbox{5cm}
\end{keepext}
\begin{keepext}
\exttask Sea water is salty, but rain is fresh. Explain how the water cycle separates the water from the salt, and name the separation technique that copies it.
\answerbox{5cm}
\end{keepext}
\end{yr8box}

% ══════════════════════════════════════════════════════════════════════════════
%  LESSON 3 — RESOURCES                                                 (Week 2 · P1)
% ══════════════════════════════════════════════════════════════════════════════
\setcounter{section}{2}
\sectiondiv{Resources}{Renewable and non-renewable resources \textbullet{} How resources are extracted \textbullet{} Benefits and risks of extraction}\label{sec:3}
\lessoncodes{VC2S8U09 \textbullet{} VC2S8H03}

\vocabhead
\begin{ctab}
\begin{tabular}{|SW{3.4cm}|W{7.0cm}|}
  \hline\hd & \thead{Word} & \thead{Definition (copy from the slide)} \\ \hline
  & \scaf{\textbf{Resource}} & \scaf{Something from the natural world that people use, such as water, wood, metals or coal.} \\ \hline
  & \textbf{Renewable} & \\ \hline
  & \textbf{Non-renewable} & \\ \hline
  & \textbf{Ore} & \\ \hline
  & \textbf{Fossil fuel} & \\ \hline
  & \textbf{Extraction} & \\ \hline
\end{tabular}
\end{ctab}

% ─── PART A ───────────────────────────────────────────────────────────────────
\lessonpart{A}{Reading and comprehension}{Teal}\label{read:3}

\begin{readingbox}{Reading --- Digging Up Australia}
Australia is one of the world's great mining countries. Iron ore from the Pilbara in Western Australia is shipped to steelworks around the world. Coal is mined in Queensland, New South Wales and in Victoria's Latrobe Valley, and natural gas is piped from under the sea. These are all \textbf{resources}: materials from the natural world that people use.

Some resources are \textbf{renewable}: they are replaced by nature in a human lifetime, like timber from a well-managed forest, or energy from the Sun and wind. Others are \textbf{non-renewable}. Coal, oil and natural gas are \textbf{fossil fuels}, formed over millions of years from the remains of ancient plants and animals. Once they are used, they are gone.

Getting a resource out of the ground is called \textbf{extraction}. Metals are found in rocks called \textbf{ores}, which are blasted and dug from mines, then crushed and processed to separate the metal. Extraction brings great benefits: jobs, money for the country, and the materials for everything from phones to bridges. It also brings risks: cleared land, lost habitat, polluted water and dust, and greenhouse gases.
\end{readingbox}

\rcreset
\rcq{Name three resources mined in Australia, and where each is found.}{2.5cm}
\rcq{What makes a resource renewable?}{2.5cm}
\rcq{Why are fossil fuels non-renewable?}{2.5cm}
\rcq{What is an ore?}{2.5cm}
\rcq{Give two benefits of extraction.}{2.5cm}
\rcq{Give two risks of extraction.}{2.5cm}

% ─── PART B ───────────────────────────────────────────────────────────────────
\lessonpart{B}{Theory and examples}{Blue}\label{theory:3}

\subsection{Renewable or Non-Renewable?}\label{sub:3.1}

\begin{skillbox}{The test}
Ask: \emph{will nature replace it within a human lifetime?} If yes, it is \textbf{renewable}. If it takes thousands or millions of years, it is \textbf{non-renewable}. A renewable resource can still run out if it is used faster than it is replaced.
\end{skillbox}

\begin{workedex}{Classifying resources}
\textbf{Timber:} renewable --- trees regrow in decades, if forests are replanted.\\
\textbf{Gold:} non-renewable --- the Earth makes no new gold in a human lifetime.\\
\textbf{Fish:} renewable, but only if not caught faster than they breed.
\end{workedex}

\begin{wedo}{Classifying resources}
Classify sunlight and natural gas as renewable or non-renewable, and give a reason for each.
\answerbox{4cm}
\end{wedo}

\subsection{Benefits and Risks of Extraction}\label{sub:3.2}

\begin{workedex}{An open-cut coal mine}
\textbf{Benefits:} jobs for the town; coal to generate electricity; money from selling coal.\\
\textbf{Risks:} cleared land and lost habitat; dust; carbon dioxide released when the coal is burned; a large hole to repair when the mine closes.
\end{workedex}

\begin{wedo}{A new lithium mine}
Lithium is used in the batteries of phones and electric cars. Give one benefit and one risk of opening a new lithium mine.
\answerbox{4cm}
\end{wedo}

\begin{countrybox}{Country and extraction}
In 2020, mining in the Pilbara destroyed the Juukan Gorge rock shelters, sacred sites of the Puutu Kunti Kurrama and Pinikura peoples with evidence of people living there for more than 46\,000 years. Deciding whether to extract a resource means weighing benefits and risks for everyone affected, including the Traditional Owners of the land.
\end{countrybox}

% ─── PART C ───────────────────────────────────────────────────────────────────
\lessonpart{C}{Practice}{Amber}\label{prac:3}

\task Sort each resource as renewable (R) or non-renewable (N), and give a reason.
\begin{ctab}
\begin{tabular}{|SW{2.8cm}|Z{1.2cm}|W{5.8cm}|}
  \hline\hd & \thead{Resource} & \thead{R or N} & \thead{Reason} \\ \hline
  & \scaf{Coal} & \scaf{N} & \scaf{Takes millions of years to form.} \\ \hline
  & Wind & & \\ \hline
  & Iron ore & & \\ \hline
  & Cotton & & \\ \hline
  & Oil & & \\ \hline
  & Fresh water & & \\ \hline
\end{tabular}
\end{ctab}

\needspace{10\baselineskip}
\task Describe the steps to get a metal from its ore, from the mine to the metal.
\answerbox{5cm}

\needspace{12\baselineskip}
\task A company wants to open a gold mine near a country town.
\pt{a}Give two benefits for the town.
\answerbox{4cm}
\pt{b}Give two risks for the environment.
\answerbox{4cm}

\needspace{10\baselineskip}
\task Fish are a renewable resource, yet some fish populations have collapsed. Explain how this can happen.
\answerbox{5cm}

% ─── PART D ───────────────────────────────────────────────────────────────────
\lessonpart{D}{Above level}{Violet}\label{ext:3}

\begin{yr8box}
These questions go beyond the Year 7 standard.

\medskip
\begin{keepext}
\exttask Recycling an aluminium can uses far less energy than making a new one from its ore. Explain why recycling metals reduces both energy use and the risks of extraction.
\answerbox{5cm}
\end{keepext}
\begin{keepext}
\exttask Electric cars use no petrol, but their batteries need mined lithium. Weigh up whether electric cars reduce the risks of extracting resources.
\answerbox{5cm}
\end{keepext}
\end{yr8box}

% ══════════════════════════════════════════════════════════════════════════════
%  LESSON 4 — RENEWABILITY AND SUSTAINABILITY                           (Week 2 · P2)
% ══════════════════════════════════════════════════════════════════════════════
\setcounter{section}{3}
\sectiondiv{Renewability and Sustainability}{What makes a use sustainable \textbullet{} Reduce, reuse, recycle, substitute \textbullet{} Melbourne's water}\label{sec:4}
\lessoncodes{VC2S8U09 \textbullet{} VC2S8H03}

\vocabhead
\begin{ctab}
\begin{tabular}{|SW{3.4cm}|W{7.0cm}|}
  \hline\hd & \thead{Word} & \thead{Definition (copy from the slide)} \\ \hline
  & \scaf{\textbf{Sustainable}} & \scaf{Using a resource in a way that leaves enough for people in the future.} \\ \hline
  & \textbf{Conserve} & \\ \hline
  & \textbf{Reduce} & \\ \hline
  & \textbf{Reuse} & \\ \hline
  & \textbf{Recycle} & \\ \hline
  & \textbf{Substitute} & \\ \hline
\end{tabular}
\end{ctab}

% ─── PART A ───────────────────────────────────────────────────────────────────
\lessonpart{A}{Reading and comprehension}{Teal}\label{read:4}

\begin{readingbox}{Reading --- When Melbourne Ran Low on Water}
Between about 1997 and 2009, south-eastern Australia went through the Millennium Drought, one of the longest dry spells on record. Less rain fell, and the rivers that fill Melbourne's dams slowed. By 2009 the dams had fallen to about a quarter full.

Melbourne had to use water more \textbf{sustainably}: in a way that would leave enough for the future. Households were asked to \textbf{conserve} water and cut their use to 155 litres per person each day. People \textbf{reduced} their use with shorter showers and water-saving taps. They \textbf{reused} water from the laundry on gardens. Some water was \textbf{recycled}, cleaned and used again on farms and parks. Rainwater tanks \textbf{substituted} for drinking water in toilets and gardens.

The state also built a desalination plant at Wonthaggi, which takes the salt out of sea water. It gives Melbourne a source of water that does not depend on rain, but it uses a great deal of energy to run.
\end{readingbox}

\rcreset
\rcq{What was the Millennium Drought?}{2.5cm}
\rcq{How full were Melbourne's dams by 2009?}{2.5cm}
\rcq{What daily target were households given?}{2.5cm}
\rcq{Give one example each of reducing, reusing and substituting water.}{3cm}
\rcq{What does a desalination plant do?}{2.5cm}
\rcq{Give one benefit and one drawback of the desalination plant.}{2.5cm}

% ─── PART B ───────────────────────────────────────────────────────────────────
\lessonpart{B}{Theory and examples}{Blue}\label{theory:4}

\subsection{What Makes a Use Sustainable?}\label{sub:4.1}

\begin{skillbox}{The sustainability test}
A use of a resource is \textbf{sustainable} if it can go on for a very long time without running the resource out or damaging the environment. For a renewable resource, use it no faster than it is replaced. For a non-renewable resource, use as little as possible and recover it where you can.
\end{skillbox}

\begin{workedex}{Is it sustainable?}
A town takes water from a river for its farms. In a dry year, it takes more water than flows into the river.

\smallskip
\textbf{Not sustainable:} it uses the water faster than the rain replaces it, so the river runs low and fish and plants downstream suffer.
\end{workedex}

\begin{wedo}{Is it sustainable?}
A school prints every worksheet single-sided and throws used paper in the general rubbish. Decide whether this is sustainable, and explain.
\answerbox{4cm}
\end{wedo}

\subsection{Four Ways to Conserve Resources}\label{sub:4.2}

\begin{skillbox}{Reduce, reuse, recycle, substitute}
\textbf{Reduce:} use less in the first place. \quad \textbf{Reuse:} use something again for the same or a new job. \quad \textbf{Recycle:} break a material down and make it into something new. \quad \textbf{Substitute:} use a different, more sustainable resource instead.
\end{skillbox}

\begin{workedex}{Plastic bottles}
\textbf{Reduce:} carry a refillable bottle. \textbf{Reuse:} refill a bottle you already have. \textbf{Recycle:} put empty bottles in the recycling bin to be made into new plastic. \textbf{Substitute:} use a steel bottle that lasts for years.
\end{workedex}

\begin{wedo}{School lunch packaging}
Give one way to reduce, reuse, recycle and substitute the packaging in a school lunch.
\answerbox{4.5cm}
\end{wedo}

% ─── PART C ───────────────────────────────────────────────────────────────────
\lessonpart{C}{Practice}{Amber}\label{prac:4}

\task Write whether each action is reduce, reuse, recycle or substitute.
\begin{ctab}
\begin{tabular}{|SW{6.8cm}|W{3.0cm}|}
  \hline\hd & \thead{Action} & \thead{Which one?} \\ \hline
  & \scaf{Turning off the tap while brushing teeth} & \scaf{Reduce} \\ \hline
  & Using a glass jar to store pencils & \\ \hline
  & Putting old newspaper in the recycling bin & \\ \hline
  & Using rainwater to flush the toilet & \\ \hline
  & Taking shorter showers & \\ \hline
\end{tabular}
\end{ctab}

\needspace{10\baselineskip}
\task Suggest three ways our school could use water more sustainably. Explain how each one helps.
\answerbox{5cm}

\needspace{10\baselineskip}
\task Is it sustainable to burn coal for electricity? Explain your answer.
\answerbox{5cm}

\needspace{10\baselineskip}
\task Explain why recycling is better than throwing things away, but reducing is better than recycling.
\answerbox{5cm}

% ─── PART D ───────────────────────────────────────────────────────────────────
\lessonpart{D}{Above level}{Violet}\label{ext:4}

\begin{yr8box}
These questions go beyond the Year 7 standard.

\medskip
\begin{keepext}
\exttask The desalination plant makes Melbourne's water supply more reliable but uses a lot of energy. Weigh up whether it makes Melbourne's water use more or less sustainable.
\answerbox{5cm}
\end{keepext}
\begin{keepext}
\exttask Fifteen people each use 155 litres of water a day. How much water do they use in a year? How many 9-litre buckets is that?
\workbox{4cm}
\end{keepext}
\end{yr8box}

% ══════════════════════════════════════════════════════════════════════════════
%  LESSON 5 — ENERGY AND RENEWABLE ENERGY                               (Week 2 · P3)
% ══════════════════════════════════════════════════════════════════════════════
\setcounter{section}{4}
\sectiondiv{Energy and Renewable Energy}{How electricity is generated \textbullet{} Fossil fuels and renewables \textbullet{} Comparing energy sources}\label{sec:5}
\lessoncodes{VC2S8U09 \textbullet{} VC2S8U15 \textbullet{} VC2S8H03}

\vocabhead
\begin{ctab}
\begin{tabular}{|SW{3.4cm}|W{7.0cm}|}
  \hline\hd & \thead{Word} & \thead{Definition (copy from the slide)} \\ \hline
  & \scaf{\textbf{Generator}} & \scaf{A machine that turns the spinning of a turbine into electricity.} \\ \hline
  & \textbf{Turbine} & \\ \hline
  & \textbf{Renewable energy} & \\ \hline
  & \textbf{Greenhouse gas} & \\ \hline
  & \textbf{Reliability} & \\ \hline
\end{tabular}
\end{ctab}

% ─── PART A ───────────────────────────────────────────────────────────────────
\lessonpart{A}{Reading and comprehension}{Teal}\label{read:5}

\begin{readingbox}{Reading --- Where Victoria's Electricity Comes From}
For most of the last century, Victoria's electricity came from brown coal dug from huge open-cut mines in the Latrobe Valley. In a coal power station, the coal is burned to boil water. The steam spins a \textbf{turbine}, and the turbine turns a \textbf{generator}, which makes electricity. Burning coal releases carbon dioxide, a \textbf{greenhouse gas} that traps heat in the atmosphere and adds to climate change.

Today more and more of Victoria's electricity comes from \textbf{renewable energy}. Wind turbines spin in the wind on hills and along the coast. Solar panels on roofs turn sunlight straight into electricity. Hydroelectric power stations in the mountains let water fall through turbines.

Each source has trade-offs. A coal power station can run day and night: its \textbf{reliability} is high. The Sun does not shine at night and the wind does not always blow, so renewable electricity has to be stored in batteries or balanced with other sources. Engineers now plan an electricity system that is clean, reliable and affordable all at once.
\end{readingbox}

\rcreset
\rcq{Where has most of Victoria's electricity come from in the past?}{2.5cm}
\rcq{Describe how a coal power station makes electricity.}{3cm}
\rcq{Why is burning coal a problem for the climate?}{2.5cm}
\rcq{Name three renewable sources of electricity.}{2.5cm}
\rcq{What does reliability mean for an energy source?}{2.5cm}
\rcq{Why does renewable electricity sometimes need to be stored?}{2.5cm}

% ─── PART B ───────────────────────────────────────────────────────────────────
\lessonpart{B}{Theory and examples}{Blue}\label{theory:5}

\subsection{How Electricity is Generated}\label{sub:5.1}

\begin{center}
\begin{tikzpicture}[font=\small, box/.style={draw=Ink, line width=0.8pt, rounded corners=1.5mm, minimum width=17mm, minimum height=9mm, align=center, fill=white}]
  \node[box] (f) at (0,0) {Fuel\\burned};
  \node[box] (w) at (2.25,0) {Water\\boiled};
  \node[box] (s) at (4.5,0) {Steam\\spins turbine};
  \node[box] (g) at (6.75,0) {Generator};
  \node[box] (e) at (9.0,0) {Electricity};
  \foreach \a/\b in {f/w,w/s,s/g,g/e}{\draw[Ink, line width=1pt, -{Stealth[length=5pt]}] (\a) -- (\b);}
\end{tikzpicture}
\figcap{In a coal or gas power station. Wind and hydro spin the turbine directly; solar panels need no turbine at all.}
\end{center}

\begin{workedex}{Energy transfers in a wind turbine}
The wind's moving energy turns the blades (\textbf{turbine}). The spinning shaft turns the \textbf{generator}, which produces \textbf{electrical} energy. No fuel is burned, so no greenhouse gases are released while it runs.
\end{workedex}

\begin{wedo}{Energy transfers in a hydro power station}
Water stored behind a dam falls through pipes to the power station. Describe how this makes electricity.
\answerbox{4cm}
\end{wedo}

\subsection{Comparing Energy Sources}\label{sub:5.2}

\begin{workedex}{Coal and solar}
\textbf{Coal:} reliable day and night; existing power stations; releases large amounts of greenhouse gas; non-renewable.\\
\textbf{Solar:} renewable; no greenhouse gas while running; panels on homes cut power bills; makes no electricity at night, so needs storage or another source.
\end{workedex}

\begin{wedo}{Wind}
Give one benefit and one drawback of wind power.
\answerbox{4cm}
\end{wedo}

% ─── PART C ───────────────────────────────────────────────────────────────────
\lessonpart{C}{Practice}{Amber}\label{prac:5}

\task Compare four energy sources.
\begin{ctab}
\begin{tabular}{|SW{1.5cm}|Z{2.1cm}|W{3.0cm}|W{3.3cm}|}
  \hline\hd & \thead{Source} & \thead{Renewable?} & \thead{Benefit} & \thead{Risk or drawback} \\ \hline
  & \scaf{Coal} & \scaf{No} & \scaf{Reliable day and night} & \scaf{Releases greenhouse gases} \\ \hline
  & Solar & & & \\ \hline
  & Wind & & & \\ \hline
  & Hydro & & & \\ \hline
  & Gas & & & \\ \hline
\end{tabular}
\end{ctab}

\needspace{10\baselineskip}
\task Explain why a solar power system often has a battery.
\answerbox{5cm}

\needspace{10\baselineskip}
\task A town must choose between a new gas power station and a wind farm with batteries. Recommend one, giving two reasons.
\answerbox{5cm}

\needspace{10\baselineskip}
\task Explain why electricity from a coal power station is not sustainable in the long term.
\answerbox{5cm}

% ─── PART D ───────────────────────────────────────────────────────────────────
\lessonpart{D}{Above level}{Violet}\label{ext:5}

\begin{yr8box}
These questions go beyond the Year 7 standard.

\medskip
\begin{keepext}
\exttask Pumped hydro stores energy by pumping water uphill when there is spare electricity, then letting it fall through turbines later. Explain how this helps a grid that uses a lot of solar and wind power.
\answerbox{5cm}
\end{keepext}
\begin{keepext}
\exttask Some people say every source of electricity has an environmental cost. Choose a renewable source and describe its environmental cost.
\answerbox{5cm}
\end{keepext}
\end{yr8box}
